\begin{proof}
\begin{enumerate}[label=\arabic*.]
\item We use the marked construction of \cref{maps:src:F.2.5} and the
integral cusp-collapse assertion of \cref{night:P1.11}.
Explicitly, $M_1=A^t$, $M_2=B^t$ and the vectors $v_1,v_2$ are
those printed in item~1 of the proof of \cref{maps:src:F.2.5};
$S$ is the displayed integral marking in item~1 of the proof of
\cref{night:P1.1}.  Write
$m_1=5$, $m_2=4$, and put
\[
\begin{split}
C_1&=\begin{pmatrix}
-1&-1&0&-1&0\\0&1&0&1&1\\0&1&0&0&0\\
0&-1&-1&0&-1\\0&0&1&0&0
\end{pmatrix},\\[2pt]
C_2&=\begin{pmatrix}
-1&0&-1&0&-1\\0&1&0&0&0\\0&0&0&1&0\\
0&-1&0&-1&1\\0&0&1&0&0
\end{pmatrix}.
\end{split}
\]
For the specified translation vectors $v_i$, the matrices
$P_i=[v_i,(0;C_i)]$ are unimodular.  Direct multiplication gives
\[
 P_i^{-1}M_iP_i=1\oplus D_i,\qquad Q_i=P_i^{-t},
\]
where $D_1$ sends $e_j$ to $e_{j+1\bmod5}$, and $D_2$ sends
$e_j$ to $e_{j+1\bmod4}$ for $0\le j<4$ and $e_4$ to $-e_4$.
These equalities give the actual integral coordinate maps.  In these
coordinates the affine action is
\[
 (u,y,z)\longmapsto
 (u+1/m_i,D_i y,e^{2\pi i/m_i}z).
\]
Thus its core $F_i$ is the mapping torus of $D_i:T^5\to T^5$.
The function $W=e^{-2\pi i u}z$ trivializes its normal disc bundle.
With $\theta=m_i u\bmod1$, the coarse-base function on that bundle is
\[
                  q=W^{m_i}e^{2\pi i\theta}.
\]
In particular the boundary is $F_i\times S^1_W$ with this specified
map to the base circle.

\item We compute the integral restriction images on this model.
Let $V=\mathbb Z^5$ be the vertical cohomology lattice, let
$T_{i,k}=\bigwedge^kD_i^t$, and let
$N_{i,k}=\sum_{a=0}^{m_i-1}T_{i,k}^{a}$.
The signed permutations preserve the product cell structure with one
vertex and one edge on each circle.  Its cellular differential is zero.
Fibre cells and their prisms give the mapping-torus complex
\[
 C^q(F_i)=\bigwedge^qV\oplus\bigwedge^{q-1}V,
 \qquad d(a,b)=(0,(T_{i,q}-I)a).
\]
The actual $m_i$-fold covering by the generic six-torus pulls back
$(a,b)$ to $(a,N_{i,q-1}b)$: the horizontal cell is the union of the
$m_i$ consecutive prisms, with their transported coefficients.
Consequently its image in $L_q=\bigwedge^q\mathbb Z^6$ is
\[
 J_{i,q}=\bigwedge^q Q_i
 \left(\operatorname{Inv}T_{i,q}\ \oplus\
 e_0\wedge N_{i,q-1}\bigwedge^{q-1}V\right).
\]
Here vertical coordinates occupy slots $1,\ldots,5$ in the six-lattice.
This formula specifies both the lattice and its inclusion.

We fix its column basis as follows.  Exterior coordinates are in
increasing lexicographic order, with indices starting at zero.
For each signed-permutation orbit, start at its least coordinate
and sum its successive signed images until it returns there.
A positive orbit of length $l$ gives a primitive invariant column;
the norm image has the column $m_i/l$ times that column.
A negative orbit has zero invariant lattice and one $\mathbb Z/2$
in its coinvariant group.  Indeed the successive relations identify
all its coordinates with signed copies of the first, and the last
relation is twice that first coordinate.  Order the positive orbits
by their least coordinates, with vertical columns before horizontal
columns.  For $D_1$ there is no torsion in vertical degrees $1$ through
$4$.  For $D_2$ there is one negative orbit in each such degree; the
positive-orbit counts are $1,2,2,1$.

\item At the cusp, \cref{night:P1.11} supplies a primitive collapse
image and free homology, with the only collapse cokernel equal to
$\mathbb Z$ in degree four.  Dualizing those integral maps therefore
gives
\[
 J_{0,q}=\operatorname{Inv}\bigl(\bigwedge^q M_0^t\bigr),
\]
with cohomological restriction kernel zero for $q=1,2,3$ and
$\mathbb Z$ for $q=4$.
Here are explicit primitive bases, which also fix the column order
in the calculation below.  In the supplied cusp marking $S$, write
the cohomology basis as $x_0,x_1,x_2,y_0,y_1,y_2$; its action sends
$x_i$ to $x_i+y_i$ and fixes $y_i$.
In degree one take $(y_0,y_1,y_2)$.  In degree two take, in order,
\[
 y_i\wedge y_j\ (i<j),\qquad
 x_i\wedge y_i\ (0\le i\le2),\qquad
 x_i\wedge y_j+x_j\wedge y_i\ (i<j).
\]
In degree three take $y_0\wedge y_1\wedge y_2$, then
$x_i\wedge y_j\wedge y_k$ in lexicographic $(j,k)$ order with
$j<k$ and $i=j,k$, and then
\[
\begin{split}
 &x_0\wedge y_1\wedge y_2+x_1\wedge y_0\wedge y_2,\\
 &x_0\wedge y_1\wedge y_2-x_2\wedge y_0\wedge y_1.
\end{split}
\]
Their coordinate equations show that these bases are primitive and
have ranks $3,9,9$.  For degree four set
$J=\left(\begin{smallmatrix}0&I_3\\-I_3&0\end{smallmatrix}\right)$,
and let $H(e_I)$ be the complementary exterior basis vector, with
the sign of the permutation $(I,I^c)$.  Apply
$H\bigwedge^2J$ to the displayed degree-two basis.
The identity $JU=U^{-t}J$ gives the invariant equation, and the
signed-permutation maps $H$ and $\bigwedge^2J$ preserve primitivity.
Transport all these columns by $\bigwedge^qS^{-t}$ to obtain
$J_{0,q}$ in the original marking.  Thus its degree-four rank is nine.

\item Let $\mathcal R_q=R^qf_*\mathbb Z$ on the coarse base $S^2$,
let $j$ include its thrice-punctured locus, and let $\mathcal S_q$
be the image of $\mathcal R_q\to j_*L_q$.
Since $M$ is compact and the base is Hausdorff, the proper-base-change
map identifies its stalks with fibre cohomology
\cite[Tag 09V5, Lemma 20.18.1; chapter PDF p.~30]{StacksProject}.
Its three stalk images are exactly the lattices above.  The kernel
$\mathcal K_q$ is supported at those points: for $q=2,3$ it is the
order-four $\mathbb Z/2$, and for $q=4$ it is that $\mathbb Z/2$
together with the cusp $\mathbb Z$.
Point-supported sheaves have no positive cohomology, so
$H^p(\mathcal R_q)=H^p(\mathcal S_q)$ for $p>0$.

Put $A_q=\bigwedge^qM_1^t$ and $B_q=\bigwedge^qM_2^t$.
The supplied based marking gives
$\bigwedge^qM_0^t=(B_qA_q)^{-1}$.
Oriented cellular duality on the punctured sphere and its two-loop
spine gives the compact-support complex $L_q^2\to L_q$ in degrees
one and two.  Attaching the three stalk lattices gives the full complex
\[
\begin{gathered}
 J_{1,q}\oplus J_{2,q}\oplus J_{0,q}
 \xrightarrow{d_0}L_q^2\xrightarrow{\Delta_q}L_q,\\
 d_0(j_1,j_2,j_0)=(j_1-j_0,j_2-A_qj_0),\\
 \Delta_q(a,b)=(A_q-I)a+(B_q-I)b.
\end{gathered}
\]
For the cusp column, the two consecutive peripheral paths carry
$j_0$ and $A_qj_0$; their boundary is
$(B_qA_q-I)j_0=0$.  The reverse cusp orientation gives their negatives.
The two finite columns are the corresponding single peripheral
loops.  This is the relative cellular sequence for extension by zero
and the three stalk inclusions, so it retains their integral indices.
It also verifies $\Delta_qd_0=0$ with the displayed transports.

Here are finite determinant certificates for its cohomology.  They
use precisely the ordered bases just given, so every matrix is
specified by the marked $M_i,S,C_i$ and exterior minors.
For each $q$ select the following columns $\mathcal P_q$ of
$\Delta_q$ and call their square matrix $E_q$:
\[
\begin{aligned}
\mathcal P_1&=(1,2,5,8,9,10),\\
\mathcal P_2&=(0,1,2,8,9,11,12,13,16,17,23,24,25,26,28),\\
\mathcal P_3&=(1,2,3,5,6,10,11,12,13,15,16,18,19,\\
             &\hspace{22mm}26,29,31,32,34,35,37),\\
\mathcal P_4&=(0,1,5,7,9,10,11,13,16,17,20,23,26,28,29).
\end{aligned}
\]
Direct integer multiplication gives
\[
 \det E_q=(2,2,2,-2)_q,\qquad
 E_q^{-1}\Delta_q\in\operatorname{Mat}(\mathbb Z).
\]
Hence $E_q\mathbb Z^{\binom6q}$ is the entire image of
$\Delta_q$ and its cokernel is $\mathbb Z/2$.
Let $\mathcal F_q$ be the increasing complement of $\mathcal P_q$
in $\{0,\ldots,2\binom6q-1\}$.
These free-coordinate rows identify $\ker\Delta_q$ integrally:
the pivot coordinates of a kernel vector are
$-E_q^{-1}(\Delta_q)_{\mathcal F_q}$ times its free coordinates.
The relation matrix is therefore $R_q=(d_0)_{\mathcal F_q}$.
Its sizes, ranks and primitive minors are
\[
\begin{array}{c|c|c|l}
q&\text{size of }R_q&\operatorname{rank}R_q&
\text{maximal minors}\\ \hline
1&6\times7&6&-20,\ -5,\ 4\\
2&15\times15&15&\det R_2=1\\
3&20\times17&17&-1\\
4&15\times15&15&\det R_4=-1
\end{array}
\]
For $q=1$ the three column sets are respectively
$(0,1,2,3,5,6)$, $(0,1,2,4,5,6)$, $(0,2,3,4,5,6)$,
using all rows.  For $q=3$ use all columns and rows
$(0,1,\ldots,13,15,16,17)$.
These determinants have greatest common divisor one, so the
quotient by the relations is free.  The stalk column counts are
$(2,2,3)$, $(3,3,9)$, $(4,4,9)$ and $(3,3,9)$.
Thus the kernel dimensions and the quotient ranks give exactly
\[
\begin{array}{c|ccc}
q&H^0(\mathcal S_q)&H^1(\mathcal S_q)&H^2(\mathcal S_q)\\ \hline
1&\mathbb Z&0&\mathbb Z/2\\
2&0&0&\mathbb Z/2\\
3&0&\mathbb Z^3&\mathbb Z/2\\
4&0&0&\mathbb Z/2.
\end{array}
\]

\item We compute the torsion differential on the actual boundary.
Let $g$ be the core-circle loop of $F_2$ and $h$ the $W$-circle loop.
Their coarse-base degrees are one and four.  The basis
\[
 v=g^4h^{-1},\qquad t=g,
 \qquad [v,t]_{[g,h]}=\begin{pmatrix}4&1\\-1&0\end{pmatrix}
\]
has determinant one; $v$ lies in the generic torus fibre and $t$
projects once around the base.
For $k=q-1$ let $T=\bigwedge^kD_2^t$ and $N=I+T+T^2+T^3$.
A torsion class $[w]\in\operatorname{coker}(T-I)$ has a mapping-torus
prism representative with values $w$ on $g$ and zero on $h$.
Subdivision into successive transported prisms gives values $Nw$
on $v$ and $w$ on $t$.  An integral equation $nw=(T-I)a$ with
$n\ne0$ implies $nNw=0$, hence $Nw=0$ in the free lattice.
Its fibre restriction is therefore zero and its first boundary
Wang component is exactly $[w]$.
The relative-base connecting map takes that component to its global
coinvariant class.  In the cellular fibre and base \v{C}ech double
complex these are precisely the two successive boundaries defining
the Leray differential.  Thus, up to the single boundary-orientation
sign,
\[
 d_2[w]=\pm[w]\in L_{q-1}/\operatorname{im}\Delta_{q-1}.
 \tag{F.6}\label{maps:eq:F.6}
\]
Changing $w$ by $(T-I)a$ changes a boundary representative by a
coboundary.  The coefficient is one because the loop $t$ has degree
one over the coarse-base circle.

For $k=1,2,3$ the negative signed orbits have representatives
$e_4$, $e_0\wedge e_2$, $e_0\wedge e_2\wedge e_4$.
Map them by the vertical inclusion followed by $\bigwedge^kQ_2$;
call the resulting columns $w_k\in L_k$.
The following rows $\rho_k$, read modulo two, annihilate
$\Delta_k$:
\[
\begin{aligned}
\rho_1&=(0,1,1,1,1,1),\\
\rho_2&=(1,1,1,1,1,1,1,0,0,1,0,1,0,1,1),\\
\rho_3&=(0,0,0,0,0,0,0,0,0,0,1,1,0,1,0,1,1,1,0,0).
\end{aligned}
\]
The explicit representatives and integer witnesses are
\[
\begin{aligned}
w_1&=(0,0,1,1,1,0)^t,\\
w_2&=(0,0,0,0,0,0,0,0,-1,0,0,-1,0,-1,-1)^t,\\
w_3&=(0,0,0,0,0,0,0,0,0,0,0,0,1,0,1,1,0,0,0,0)^t,\\
z_1&=(0,0,0,-1,-1,-1)^t,\\
z_2&=(-1,0,1,-2,-2,-2,4,2,-1,-3,2,0,-2,0,2)^t,\\
z_3&=(-2,-2,-3,0,0,-1,-1,-1,2,2,-1,3,-1,4,3,\\
   &\hspace{23mm}-3,-2,1,-4,1)^t.
\end{aligned}
\]
They satisfy $E_kz_k=2w_k$ and $\rho_kw_k=1\bmod2$.
The nonzero local torsion class therefore maps to the nonzero global
class in each degree $q=2,3,4$.

\item Since $H^0(\mathcal S_q)=0$ in these three degrees, the
point-supported kernels give
\[
 H^0(\mathcal R_2)=\mathbb Z/2,\quad
 H^0(\mathcal R_3)=\mathbb Z/2,\quad
 H^0(\mathcal R_4)=\mathbb Z\oplus\mathbb Z/2.
\]
The Leray sequence
\cite[Tag 01F2, Lemma 20.13.4; chapter PDF p.~22]{StacksProject}
has only base columns zero, one and two.
The maps from $(0,2)$ to $(2,1)$ and from $(0,3)$ to $(2,2)$ are
the isomorphisms just proved.  The restriction of the map from
$(0,4)$ to $(2,3)$ to its torsion summand is also an isomorphism.
Writing the image of a free cusp generator as $\epsilon\in\mathbb Z/2$,
its kernel is the free group generated by $(1,\epsilon)$.
All total-degree-three terms consequently vanish.  In total degree
four the remaining terms are $\mathbb Z$ and $\mathbb Z^3$.
Their extension is free, so
\[
                   H^3(M;\mathbb Z)=0,\qquad
                   H^4(M;\mathbb Z)=\mathbb Z^4.
\]
By \cref{maps:src:F.2.5}, $\pi_1(M)=0$ and the rational homology
in degrees two and three vanishes.  Universal coefficients
\cite[Theorem 3.2, p.~195, and Corollary 3.3, p.~196]{Hatcher2002}
identify the torsion of $H_2$ and $H_3$ with the corresponding
$\operatorname{Ext}$ subgroups of $H^3$ and $H^4$; both are zero.
These finitely generated groups have rational rank zero, hence
$H_2=H_3=0$.  Integral Poincar\'e duality
\cite[Theorem 3.30, p.~241]{Hatcher2002} determines the remaining
groups as $\mathbb Z,\mathbb Z^4,\mathbb Z$ in degrees $0,4,8$ and
zero otherwise.  Finally, the Hurewicz homomorphisms
\cite[Theorem 4.32, p.~366]{Hatcher2002} give successively
$\pi_2=H_2=0$ and $\pi_3=H_3=0$.
\end{enumerate}
\end{proof}
